\chapter{Wstęp}

\section{Motywacja}

Poczta elektroniczna pozostaje jednym z podstawowych kanałów komunikacji
prywatnej i zawodowej. Jej powszechność pozwala niewielkim kosztem rozsyłać
wiadomości reklamowe, phishingowe i oszukańcze do dużej liczby odbiorców.
Filtr musi rozpoznawać zmieniające się sposoby formułowania wiadomości i nie
blokować poprawnej korespondencji. Fałszywy alarm może oznaczać przeoczenie
ważnej wiadomości, natomiast przepuszczenie ataku naraża odbiorcę na utratę
danych lub pieniędzy.

Jednym z etapów wielowarstwowego filtrowania poczty jest klasyfikacja tekstu.
W klasycznych metodach treść reprezentuje się między innymi za pomocą TF-IDF,
a następnie przekazuje do klasyfikatora. Modele językowe wykorzystują natomiast
reprezentacje zależne od kontekstu oraz wiedzę zdobytą podczas treningu
wstępnego. Ich zastosowanie wiąże się jednak z wyższym kosztem treningu i
inferencji niż w przypadku prostszych klasyfikatorów. Filtr pocztowy może
przetwarzać duży strumień wiadomości, dlatego oprócz skuteczności znaczenie mają
również czas odpowiedzi, przepustowość i zużycie zasobów. Powstaje zatem
pytanie, czy modele językowe uzyskują lepsze wyniki w klasyfikacji oraz jaki
jest ich koszt w porównaniu z klasycznymi klasyfikatorami.

\section{Cel pracy}

Celem pracy jest sprawdzenie, czy mały model językowy dostrojony za pomocą
adapterów LoRA może skutecznie klasyfikować wiadomości elektroniczne jako
\texttt{ham} albo \texttt{spam}. Badanie ma określić wpływ parametrów treningu
i sposobu sformułowania zadania na wyniki uzyskiwane na kilku korpusach
wiadomości.

Drugim celem jest porównanie dostrojonego modelu z modelem bez dostrajania oraz
prostszymi klasyfikatorami tekstu. Porównanie obejmuje skuteczność, profil
błędów i koszt inferencji.

\section{Zakres i metoda badawcza}

Badania dotyczą klasyfikacji wiadomości na podstawie tematu i treści. Nie
analizowano załączników, reputacji adresów i domen, pełnych nagłówków
transportowych ani wyników SPF, DKIM i DMARC. Przyjęta klasa \texttt{spam}
obejmuje niechciane reklamy, phishing oraz wiadomości oszukańcze. Model
rozstrzyga, czy wiadomość jest pożądana, ale nie określa rodzaju wykrytego
nadużycia.

Jako model bazowy wybrano \texttt{Qwen3-0.6B}, który dostrajano metodą LoRA.
Dla klasyfikacji przez ocenę następnego tokenu porównano 16 konfiguracji
różniących się długością kontekstu, współczynnikiem uczenia i parametrami
adaptera. Oceniono również pośrednie punkty kontrolne. Wybrane ustawienia
wykorzystano następnie do porównania czterech metod: klasyfikacji sekwencji,
generowania etykiety, oceny następnego tokenu oraz generowania odpowiedzi
ustrukturyzowanej.

Wyniki odniesiono do modelu bez dostrajania. Porównanie objęło Naive Bayes,
regresję logistyczną i SVM korzystające z reprezentacji TF-IDF, a także
\texttt{fastText}, MiniLM z regresją logistyczną oraz DistilBERT. Dla metod
bazowych dobrano parametry i progi klasyfikacji na wydzielonym zbiorze
walidacyjnym. Końcowe porównanie przeprowadzono na czterech zbiorach testowych,
po 2000 wiadomości w każdym. Koszt inferencji zmierzono na wspólnej próbce 1000
wiadomości.

\section{Pytania badawcze}

Zakres pracy opisują następujące pytania badawcze:

\begin{enumerate}
    \item Jakie wyniki osiąga mały model językowy dostrojony za pomocą LoRA na
    próbce treningowej, w zewnętrznej walidacji i w teście końcowym?
    \item Który z badanych sposobów użycia modelu językowego daje
    najkorzystniejsze połączenie skuteczności, niezależności decyzji od formatu
    generowanej odpowiedzi i prostoty inferencji?
    \item Jak parametry treningu oraz jego czas trwania wpływają na wynik
    modelu na korpusach niewykorzystanych w treningu?
    \item Jaki kompromis między skutecznością a kosztem inferencji wynika z
    porównania dostrojonego modelu językowego z prostszymi klasyfikatorami?
\end{enumerate}

\section{Struktura pracy}

Po wstępie przedstawiono podstawy teoretyczne klasyfikacji wiadomości oraz
działania modeli językowych. Rozdział obejmuje tokenizację, architekturę
Transformer, dostrajanie adapterowe i sposoby wykorzystania modelu językowego
do klasyfikacji. Następnie omówiono mechanizmy stosowane w filtrowaniu poczty,
od reguł i reputacji po klasyczne algorytmy uczenia maszynowego oraz
klasyfikatory transformerowe.

Część badawcza opisuje wybór modelu, przygotowanie zbiorów danych i
przebieg eksperymentów. Kolejne sekcje dotyczą doboru parametrów LoRA, analizy
punktów kontrolnych, porównania sposobów dostrajania oraz zestawienia
wariantów Qwen3 z metodami bazowymi i modelem bez dostrajania na wspólnym
teście końcowym. Rozdział końcowy odpowiada na pytania badawcze, omawia
ograniczenia, możliwość wdrożenia klasyfikatora i dalsze kierunki badań.
